\documentclass[11pt]{article}
\usepackage[a4paper,margin=27mm]{geometry}
\usepackage{amsmath,amssymb,mathtools}
\usepackage{fontspec,unicode-math}
\setmainfont{Latin Modern Roman}
\setmathfont{Latin Modern Math}
\usepackage{xurl}
\usepackage[hidelinks]{hyperref}
\newcommand{\sref}[2]{\hyperlink{source:#1:#2}{[S#2]}}
\newcommand{\cataloguescope}[1]{\paragraph{Mathematical question.}#1}
\hypersetup{pdftitle={UnsolvedMath Problem 8900010}}
\begin{document}
% UnsolvedMath ID 8900010; source catalogue code AMR-088-0010
\setcounter{section}{8900009}
\section{L-value criterion for congruent numbers}\label{8900010}
{\small\sffamily UnsolvedMath ID 8900010 · Graph Theory}
\subsection{Definitions and mathematical statement}

\paragraph{Shared notation.}$\mathbb N=\{1,2,\ldots\}$, and $\mathbb Z,\mathbb Q,\mathbb R,\mathbb C$ denote the integers, rationals, reals and complex numbers. Any other conventions are specified locally.


A positive rational number $q$ is \emph{congruent} if there are positive rational numbers $a,b,c$ with $a^2+b^2=c^2$ and $ab/2=q$; thus $q$ is the area of a right triangle with rational side lengths.
For $n\in\mathbb N$ let $E_n/\mathbb Q$ be the elliptic curve $y^2=x^3-n^2x$, including its point at infinity. Its Hasse--Weil $L$-function is defined initially for $\operatorname{Re}s>3/2$ by
\[
L(E_n,s)=\prod_{p\text{ good}}(1-a_pp^{-s}+p^{1-2s})^{-1}
\prod_{p\text{ bad}}(1-a_pp^{-s})^{-1}.
\]
Here $p$ ranges over primes, good or bad according to the minimal model of $E_n$; at a good prime $a_p=p+1-\#E_n(\mathbb F_p)$, and at a bad prime $a_p$ is $1$, $-1$ or $0$ according as reduction is split multiplicative, nonsplit multiplicative or additive. The value $L(E_n,1)$ is taken using analytic continuation. Omitting finitely many Euler factors as in the slides does not change whether this value is zero, since those factors are finite and nonzero at $s=1$.
\paragraph{Statement recorded in the supplied source.}
For every $n\in\mathbb N$, is $n$ congruent if and only if $L(E_n,1)=0$? \sref{8900010}{2}
\subsection{Short English statement}
For each congruent-number elliptic curve, does vanishing of its $L$-function at one exactly detect the existence of a rational right triangle of area $n$?
\subsection{Sources}
\begin{description}
\item[\hypertarget{source:8900010:1}{S1}] ULAM.AI and the UnsolvedMath Contributors, UnsolvedMath: A Curated Collection of Open Mathematics Problems, record 8900010 (AMR-088-0010). CC BY 4.0. \url{https://huggingface.co/datasets/ulamai/UnsolvedMath}.
\item[\hypertarget{source:8900010:2}{S2}] Christian Wuthrich, \emph{Open problems in number theory}, lecture slides (December 2011) slides headed Conjecture (part of Birch and Swinnerton-Dyer conjecture). \url{https://www.maths.nottingham.ac.uk/plp/pmzcw/download/open.pdf}.
\end{description}
\label{end:8900010}
\end{document}
